\section{Experiments}





\subsection{Datasets}

\begin{table}
    \centering
    \caption{Statistics of datasets.}
    \label{tab:dataset}
    \begin{adjustbox}{max width=\linewidth}
    \begin{tabular}{l r r r r r}
    \toprule
    Datasets & \#users & \#items & \#actions & Avg. length &  Density\\
    \midrule
    Beauty & \num{40226} & \num{54542} & 0.35m & 8.8 & 0.02\% \\
    Steam & \num{281428} & \num{13044} & 3.5m & 12.4 & 0.10\% \\
    ML-1m & \num{6040} & \num{3416} & 1.0m & 163.5 & 4.79\% \\
    ML-20m & \num{138493} & \num{26744} & 20m & 144.4 & 0.54\%\\
    \bottomrule
    \end{tabular}
    \end{adjustbox}
\end{table}

We evaluate the proposed model on four real-world representative  datasets which vary significantly in domains and sparsity.

\begin{itemize}
    \item Amazon \textbf{Beauty}\footnote{\url{http://jmcauley.ucsd.edu/data/amazon/}}: This is a series of product review datasets crawled from Amazon.com by \citeauthor{McAuley:sigir2015}~\cite{McAuley:sigir2015}. 
    They split the data into separate datasets according to the top-level product categories on Amazon.
    In this work, we adopt the ``Beauty'' category.
    \item \textbf{Steam}\footnote{\url{https://cseweb.ucsd.edu/~jmcauley/datasets.html\#steam_data}}: 
    This is a dataset collected from Steam, a large online video game distribution platform, by \citeauthor{Kang:ICDM2018:SAS}~\cite{Kang:ICDM2018:SAS}.
    \item \textbf{MovieLens}~\cite{Harper:2015:MDH}: This is a popular benchmark dataset for evaluating recommendation algorithms.
    In this work, we adopt two well-established versions, MovieLens 1m (\textbf{ML-1m})\footnote{\url{https://grouplens.org/datasets/movielens/1m/}} and MovieLens 20m (\textbf{ML-20m})\footnote{\url{https://grouplens.org/datasets/movielens/20m/}}.
\end{itemize}

For dataset preprocessing, we follow the common practice in \cite{Rendle:WWW2010:FPM,Tang:WSDM2018:caser,Kang:ICDM2018:SAS}.
For all datasets, we convert all numeric ratings or the presence of a review to implicit feedback of 1 (\textit{i.e.}, the user interacted with the item).
After that, we group the interaction records by users and build the interaction sequence for each user by sorting these interaction records according to the timestamps.
To ensure the quality of the dataset, following the common practice~\cite{Rendle:WWW2010:FPM,He:www2017:NCF,Tang:WSDM2018:caser,Kang:ICDM2018:SAS}, we keep users with at least five feedbacks.
The statistics of the processed datasets are summarized in Table~\ref{tab:dataset}.

\subsection{Task Settings \& Evaluation Metrics}

To evaluate the sequential recommendation models, we adopted the \textit{leave-one-out} evaluation (\textit{i.e.}, next item recommendation) task, which has been widely used in~\cite{He:www2017:NCF,Tang:WSDM2018:caser,Kang:ICDM2018:SAS}.
For each user, we hold out the last item of the behavior sequence as the test data, treat the item just before the last as the validation set, and utilize the remaining items for training.
For easy and fair evaluation, we follow the common strategy in \cite{He:www2017:NCF,Tang:WSDM2018:caser,Kang:ICDM2018:SAS}, pairing each ground truth item in the test set with 100 randomly sampled \textit{negative} items that the user has not interacted with.
To make the sampling reliable and representative~\cite{Huang:sigir2018:ISR}, these 100 negative items are sampled according to their popularity.
Hence, the task becomes to rank these negative items with the ground truth item for each user.

\textbf{Evaluation Metrics}.
To evaluate the ranking list of all the models, we employ a variety of evaluation metrics, including \textit{Hit Ratio} (HR), \textit{Normalized Discounted Cumulative Gain} (NDCG), and \textit{Mean Reciprocal Rank} (MRR).
Considering we only have one ground truth item for each user, HR@$k$ is equivalent to Recall@$k$ and proportional to Precision@$k$; MRR is equivalent to Mean Average Precision (MAP).
In this work, we report HR and NDCG with $k={1,5,10}$.
For all these metrics, the higher the value, the better the performance.

\iffalse
$\bm{\mathrm{HR}@k}$ measures the proportion of cases that the desired item is among the top-$k$ items in all test cases.
It is computed as: 
\begin{equation*}
    \mathrm{HR}@k = \frac{1}{\vert \mathcal{U} \vert}\sum_{u \in \mathcal{U}} \mathbbm{1}(\mathrm{R}_{u, g_u} \leq k)
\end{equation*}
where, $g_u$ is the ground truth item for user $u$, $\mathrm{R}_{u, g_u}$ is the rank generated by the model for item $g_u$ and user $u$, and $\mathbbm{1}$ is an indicator function.
$\mathrm{HR}@k$ does not consider the actual rank of $g_u$ as long as it is among the top-$k$.

$\bm{\mathrm{NDCG}@k}$ is a position-aware metric which assigns larger weights on higher ranks. For next item recommendation, it is calculated as :
\begin{equation*}
    \mathrm{NDCG}@k = \frac{1}{\vert \mathcal{U} \vert}\sum_{u \in \mathcal{U}} \frac{2^{\mathbbm{1}(\mathrm{R}_{u, g_u} \leq k)} - 1}{\log_2 (\mathrm{R}_{u, g_u} + 1)}
\end{equation*}

\textbf{MRR} is the average of reciprocal ranks of the ground truth $\mathrm{R}_{u, g_u}$:
\begin{equation*}
    \mathrm{MRR} = \frac{1}{\vert \mathcal{U} \vert}\sum_{u \in \mathcal{U}}\frac{1}{\mathrm{R}_{u, g_u}}
\end{equation*}
\fi


\subsection{Baselines \& Implementation Details}

To verify the effectiveness of our method, we compare it with the following representative baselines:
\begin{itemize}[leftmargin=0.6cm]
    \item \textbf{POP}: It is the simplest baseline that ranks items according to their popularity judged by the number of interactions.
    \item \textbf{BPR-MF}~\cite{Rendle:UAI2009:BBP}: It optimizes the matrix factorization with implicit feedback using a pairwise ranking loss.
    \item \textbf{NCF}~\cite{He:www2017:NCF}: It models user–item interactions with a MLP instead of the inner product in matrix factorization.
    \item \textbf{FPMC}~\cite{Rendle:WWW2010:FPM}: It captures users' general taste as well as their sequential behaviors by combing MF with first-order MCs.
    \item \textbf{GRU4Rec}~\cite{Hidasi:ICLR2016:gru4rec}: It uses GRU with ranking based loss to model user sequences for session based recommendation.
    \item \textbf{GRU4Rec$^{+}$}~\cite{Hidasi:CIKM2018:RNN}: It is an improved version of GRU4Rec with a new class of loss functions and sampling strategy.
    \item \textbf{Caser}~\cite{Tang:WSDM2018:caser}: It employs CNN in both horizontal and vertical way to model high-order MCs for sequential recommendation.
    \item \textbf{SASRec}~\cite{Kang:ICDM2018:SAS}: It uses a left-to-right Transformer language model to capture users' sequential behaviors, and achieves state-of-the-art performance on sequential recommendation.
\end{itemize}


For NCF\footnote{\url{https://github.com/hexiangnan/neural_collaborative_filtering}}, GRU4Rec\footnote{\label{gru4rec}\url{https://github.com/hidasib/GRU4Rec}}, GRU4Rec$^+$\textsuperscript{\ref{gru4rec}}, Caser\footnote{\url{https://github.com/graytowne/caser_pytorch}}, and SASRec\footnote{\url{https://github.com/kang205/SASRec}}, we use code provided by the corresponding authors.
For BPR-MF and FPMC, we implement them using \texttt{TensorFlow}.
For common hyper-parameters in all models, we consider the hidden dimension size $d$ from $\{16, 32, 64, 128, 256\}$, the $\ell_2$ regularizer from \{1, 0.1, 0.01, 0.001, 0.0001\}, and dropout rate from \{$0, 0.1, 0.2,\cdots, 0.9$\}.
All other hyper-parameters (\textit{e.g.}, Markov order in Caser) and initialization strategies are either followed the suggestion from the methods' authors or tuned on the validation sets.
We report the results of each baseline under its optimal hyper-parameter settings.

We implement BERT4Rec\footnote{\url{https://github.com/FeiSun/BERT4Rec}} with \texttt{TensorFlow}.
All parameters are initialized using truncated normal distribution in the range $[-0.02, 0.02]$.
We train the model using Adam~\cite{Kingma:iclr2015:Adam} with learning rate of 1e-4, $\beta_1 = 0.9, \beta_2 = 0.999$, $\ell_2$ weight decay of 0.01, and linear decay of the learning rate.
The gradient is clipped when its $\ell_2$ norm exceeds a threshold of 5. %
For fair comparison, we set the layer number $L=2$ and head number $h=2$ and use the same maximum sequence length as in~\cite{Kang:ICDM2018:SAS}, $N=200$ for ML-1m and ML-20m, $N=50$ for Beauty and Steam datasets. %
For head setting, we empirically set the dimensionality of each head as 32 (single head if $d<32$).
We tune the mask proportion $\rho$ using the validation set, resulting in $\rho=0.6$ for Beauty, $\rho=0.4$ for Steam, and $\rho=0.2$ for ML-1m and ML-20m. %
All the models are trained from scratch without any pre-training on a single NVIDIA GeForce GTX 1080 Ti GPU with a batch size of 256.

\subsection{Overall Performance Comparison}
\begin{table*}[]
\setlength{\tabcolsep}{0.62em}
    \centering
    \caption{Performance comparison of different methods on next-item prediction. Bold scores are the best in each row, while underlined scores are the second best. Improvements over baselines are statistically significant with $p<0.01$.} %
    \begin{adjustbox}{max width=\textwidth}
        \begin{tabular}{l l c c c c c c c c c c}
        \toprule
        Datasets & Metric & POP & BPR-MF & NCF & FPMC & GRU4Rec & GRU4Rec$^+$ & Caser & SASRec & BERT4Rec & Improv.\\
        \midrule
        \multirow{6}{*}{Beauty} & HR@1 & 0.0077 & 0.0415 & 0.0407 & 0.0435 & 0.0402 & 0.0551 & 0.0475 & \underline{0.0906} & \textbf{0.0953} & ~~5.19\% \\
         & HR@5 & 0.0392 & 0.1209 & 0.1305 & 0.1387 & 0.1315 & 0.1781 & 0.1625 & \underline{0.1934} & \textbf{0.2207} & 14.12\%\\
         & HR@10 & 0.0762 & 0.1992 & 0.2142 & 0.2401 & 0.2343 & 0.2654 & 0.2590 & \underline{0.2653} & \textbf{0.3025} & 14.02\%\\
         & NDCG@5 & 0.0230 & 0.0814 & 0.0855 & 0.0902 & 0.0812 & 0.1172 & 0.1050 & \underline{0.1436} & \textbf{0.1599} & 11.35\%\\
         & NDCG@10 & 0.0349 & 0.1064 & 0.1124 & 0.1211 & 0.1074 & 0.1453 & 0.1360 & \underline{0.1633} & \textbf{0.1862} & 14.02\%\\
          & MRR & 0.0437 & 0.1006 & 0.1043 & 0.1056 & 0.1023 & 0.1299 & 0.1205 & \underline{0.1536} & \textbf{0.1701} & 10.74\%\\
         \midrule
        \multirow{6}{*}{Steam} & HR@1 & 0.0159 & 0.0314 & 0.0246 & 0.0358 & 0.0574 & 0.0812 & 0.0495 & \underline{0.0885} & \textbf{0.0957} & ~~8.14\%\\
         & HR@5 & 0.0805 & 0.1177 & 0.1203 & 0.1517 & 0.2171 & 0.2391 & 0.1766 & \underline{0.2559} & \textbf{0.2710} & ~~5.90\%\\
         & HR@10 & 0.1389 & 0.1993 & 0.2169 & 0.2551 & 0.3313 & 0.3594 & 0.2870 & \underline{0.3783} & \textbf{0.4013} & ~~6.08\% \\
         & NDCG@5 & 0.0477 & 0.0744 & 0.0717 & 0.0945 & 0.1370 & 0.1613 & 0.1131 & \underline{0.1727} & \textbf{0.1842} & ~~6.66\% \\
         & NDCG@10 & 0.0665 & 0.1005 & 0.1026 & 0.1283 & 0.1802 & 0.2053 & 0.1484 & \underline{0.2147} & \textbf{0.2261} & ~~5.31\%\\
         & MRR & 0.0669 & 0.0942 & 0.0932 & 0.1139 & 0.1420 & 0.1757 & 0.1305 & \underline{0.1874} & \textbf{0.1949} & ~~4.00\%\\
         \midrule
        \multirow{6}{*}{ML-1m} & HR@1 & 0.0141 & 0.0914 & 0.0397 & 0.1386 & 0.1583 & 0.2092 & 0.2194 & \underline{0.2351} & \textbf{0.2863} & 21.78\%\\
         & HR@5 & 0.0715 & 0.2866 & 0.1932 & 0.4297 & 0.4673 & 0.5103 & 0.5353 & \underline{0.5434} & \textbf{0.5876} & ~~8.13\%\\
         & HR@10 & 0.1358 & 0.4301 & 0.3477 & 0.5946 & 0.6207 & 0.6351 &  \underline{0.6692} & 0.6629 & \textbf{0.6970} & ~~4.15\%\\
         & NDCG@5 & 0.0416 & 0.1903 & 0.1146 & 0.2885 & 0.3196 & 0.3705 & 0.3832 & \underline{0.3980} & \textbf{0.4454} & 11.91\%\\
         & NDCG@10 & 0.0621 & 0.2365 & 0.1640 & 0.3439 & 0.3627 & 	0.4064 & 0.4268 & \underline{0.4368} & \textbf{0.4818} & 10.32\%\\
         & MRR & 0.0627 & 0.2009 & 0.1358 & 0.2891 & 0.3041 & 	0.3462 & 0.3648 & \underline{0.3790} & \textbf{0.4254} & 12.24\%\\
         \midrule
        \multirow{6}{*}{ML-20m} & HR@1 & 0.0221 & 0.0553 & 0.0231 & 0.1079 & 0.1459 & 0.2021 & 0.1232 & \underline{0.2544} & \textbf{0.3440} & 35.22\%\\
         & HR@5 & 0.0805 & 0.2128 & 0.1358 & 0.3601 & 0.4657 & 0.5118 & 0.3804 & \underline{0.5727} & \textbf{0.6323} & 10.41\%\\
         & HR@10 & 0.1378 & 0.3538 & 0.2922 & 0.5201 & 0.5844 & 0.6524 & 0.5427 & \underline{0.7136} & \textbf{0.7473} & ~~4.72\% \\
         & NDCG@5 & 0.0511 & 0.1332 & 0.0771 & 0.2239 & 0.3090 & 0.3630 & 0.2538 & \underline{0.4208} & \textbf{0.4967} & 18.04\%\\
         & NDCG@10 & 0.0695 & 0.1786 & 0.1271 & 0.2895 & 0.3637 & 	0.4087 & 0.3062 & \underline{0.4665} & \textbf{0.5340} & 14.47\% \\
         & MRR & 0.0709 & 0.1503 & 0.1072 & 0.2273 & 0.2967 & 0.3476 & 0.2529 & \underline{0.4026} & \textbf{0.4785} & 18.85\% \\
        \bottomrule
        \end{tabular}
    \end{adjustbox}
    \label{tab:result}
\end{table*}


Table~\ref{tab:result} summarized the best results of all models on four benchmark datasets.
The last column is the improvements of BERT4Rec relative to the best baseline.
We omit the NDCG@1 results since it is equal to HR@1 in our experiments.
It can be observed that:

The non-personalized POP method gives the worst performance\footnote{What needs to be pointed out is that such low scores for POP are because the negative samples are sampled according to the items' popularity.} on all datasets since it does not model user's personalized preference using the historical records.
Among all the baseline methods, sequential methods (\textit{e.g.}, FPMC and GRU4Rec$^+$) outperforms non-sequential methods (\textit{e.g.}, BPR-MF and NCF) on all datasets consistently.
Compared with BPR-MF, the main improvement of FPMC is that it models users' historical records in a sequential way.
This observation verifies that considering sequential information is beneficial for improving performances in recommendation systems.

Among sequential recommendation baselines, Caser outperforms FPMC on all datasets especially for the dense dataset ML-1m, suggesting that high-order MCs is beneficial for sequential recommendation.
However, high-order MCs usually use very small order $L$ since they do not scale well with the order $L$.
This causes Caser to perform worse than GRU4Rec$^+$ and SASRec, especially on sparse datasets.
Furthermore, SASRec performs distinctly better than GRU4Rec and GRU4Rec$^+$, suggesting that self-attention mechanism is a more powerful tool for sequential recommendation.


According to the results, it is obvious that BERT4Rec performs best  among all methods on four datasets in terms of all evaluation metrics.
It gains \textbf{7.24\%} HR@10, \textbf{11.03\%} NDCG@10, and \textbf{11.46\%} MRR improvements (on average) against the strongest baselines.

\noindent \textbf{Question 1}: \textit{Do the gains come from the bidirectional self-attention model or from the Cloze objective?}

To answer this question, we try to isolate the effects of these two factors by constraining the Cloze task to mask only one item at a time.
In this way, the main difference between our BERT4Rec (with 1 mask) and SASRec is that BERT4Rec predicts the target item jointly conditioning on both left and right context.
We report the results on Beauty and ML-1m with $d=256$ in Table~\ref{tab:analysis} due to the space limitation.
The results show that BERT4Rec with 1 mask significantly outperforms SASRec on all metrics.
It demonstrates the importance of bidirectional representations for sequential recommendation.
Besides, the last two rows indicate that the Cloze objective also improves the performances.
Detailed analysis of the mask proportion $\rho$ in Cloze task can be found in \S~\ref{sec:mask}

\begin{table}
    \centering
    \caption{Analysis on bidirection and Cloze with $d=256$.}
    \label{tab:analysis}
    \begin{adjustbox}{max width=\linewidth}
    \begin{tabular}
        {l c c c c c c}
        \toprule
       \multirow{2}{*}{Model} & \multicolumn{3}{c}{Beauty} & \multicolumn{3}{c}{ML-1m} \\ \cmidrule(lr){2-4} \cmidrule(lr){5-7}
        & HR@10 & NDCG@10 & MRR & HR@10 & NDCG@10 & MRR \\ \midrule
        SASRec & 0.2653 & 0.1633 & 0.1536 & 0.6629 & 0.4368 & 0.3790 \\
        BERT4Rec (1 mask) & 0.2940 & 0.1769 & 0.1618 & 0.6869 & 0.4696 & 0.4127 \\
        BERT4Rec & 0.3025 & 0.1862 & 0.1701 & 0.6970 & 0.4818 & 0.4254 \\ 
        \bottomrule
    \end{tabular}
    \end{adjustbox}
\end{table}



\noindent \textbf{Question 2}: \textit{Why and how does bidirectional model outperform unidirectional models?}
\begin{figure}
\vspace{2mm}
    \centering
    \begin{subfigure}[b]{0.45\linewidth}
        \includegraphics[width=\textwidth]{./figs/256_BERT_Attention_layer_0_head_1}
        \caption{Layer 1, head 1}
        \label{fig:l1h1}
    \end{subfigure}
    \quad
    \begin{subfigure}[b]{0.45\linewidth}
        \includegraphics[width=\textwidth]{./figs/256_BERT_Attention_layer_0_head_2}
        \caption{Layer 1, head 2}
        \label{fig:l1h2}
    \end{subfigure}
    \begin{subfigure}[b]{0.47\linewidth}
     ~~~\includegraphics[width=\textwidth]{./figs/256_BERT_Attention_layer_1_head_2}
        \caption{Layer 2, head 2}
        \label{fig:mouse}
    \end{subfigure}
    \quad
    \begin{subfigure}[b]{0.47\linewidth}
        \includegraphics[width=\textwidth]{./figs/256_BERT_Attention_layer_1_head_4}
        \caption{Layer 2, head 4}
        \label{fig:mouse}
    \end{subfigure}
    \caption{Heat-maps of average attention weights on Beauty, the last position ``9'' denotes ``\texttt{[mask]}'' (best viewed in color).}
    \label{fig:att_vis}
\end{figure}

To answer this question, we try to reveal meaningful patterns by visualizing the average attention weights of the last 10 items during the test on Beauty in Figure~\ref{fig:att_vis}. 
Due to the space limitation, we only report four representative attention heat-maps in different layers and heads.

\input{dim_exp_fig}
We make several observations from the results.
\begin {enumerate*}[label=\bfseries\itshape\alph*\upshape)]
\item Attention varies across different heads.
For example, in layer 1, head 1 tends to attend on items at the left side while head 2 prefers to attend on items on the right side.
\item Attention varies across different layers.
Apparently, attentions in layer 2 tend to focus on more recent items.
This is because layer 2 is directly connected to the output layer and the recent items play a more important role in predicting the future.
Another interesting pattern is that heads in Figure~\ref{fig:l1h1} and~\ref{fig:l1h2} also tend to attend on \texttt{[mask]}\footnote{This phenomenon also exists in text sequence modeling using BERT.}.
It may be a way for self-attention to propagate sequence-level state to the item level.
\item Finally and most importantly, unlike unidirectional model can only attend on items at the left side, items in BERT4Rec tend to attend on the items at both sides.
This indicates that bidirectional is essential and beneficial for user behavior sequence modeling.
\end {enumerate*}



In the following studies, we examine the impact of the hyper-parameters, including the hidden dimensionality $d$, the mask proportion $\rho$, and the maximum sequence length $N$.
We analyze one hyper-parameter at a time by fixing the remaining hyper-parameters at their optimal settings. 
Due to space limitation, we only report NDCG@10 and HR@10 for the follow-up experiments. %



\subsection{Impact of Hidden Dimensionality $d$}

We now study how the hidden dimensionality $d$ affects the recommendation performance.
Figure~\ref{fig:dim_exp} shows NDCG@10 and HR@10 for neural sequential methods with the hidden dimensionality $d$ varying from 16 to 256 while keeping other optimal hyper-parameters unchanged.
We make some observations from this figure.




The most obvious observation from these sub-figures is that the performance of each model tends to converge as the dimensionality increases.
A larger hidden dimensionality does not necessarily lead to better model performance, especially on sparse datasets like Beauty and Steam.
This is probably caused by overfitting.
In terms of details, Caser performs unstably on four datasets, which might limit its usefulness.
Self-attention based methods (\textit{i.e.}, SASRec and BERT4Rec) achieve superior performances on all datasets. %
Finally, our model consistently outperforms all other baselines on all datasets even with a relatively small hidden dimensionality.
Considering that our model achieves satisfactory performance with $d {\geq} 64$, we only report the results with $d{=}64$ in the following analysis.

\subsection{Impact of Mask Proportion $\rho$}
\label{sec:mask}

\pgfplotsset{
axis background/.style={fill=gallery},
grid=both,
  xtick pos=left,
  ytick pos=left,
  tick style={
    major grid style={style=white,line width=1pt},
    minor grid style=bgc,
    draw=none
    },
  minor tick num=1,
  ymajorgrids,
	major grid style={draw=white},
	y axis line style={opacity=0},
	tickwidth=0pt,
}


\begin{figure}[ht]
\centering
    \begin{tikzpicture}[scale=0.5]
	\begin{groupplot}[
	    group style={group size=2 by 1,
	        horizontal sep = 48pt}, 
	    xlabel=\Large Dimensionality,
        ylabel=\large HR@10,
        xticklabels={0.1, 0.2, 0.3, 0.4, 0.5, 0.6, 0.7, 0.8, 0.9},
        xtick={1,2,3,4,5,6,7,8,9},
        ymajorgrids,
        major grid style={draw=white},
        y axis line style={opacity=0},
        tickwidth=0pt,
        yticklabel style={
        /pgf/number format/fixed,
        /pgf/number format/precision=5
        },
        scaled y ticks=false,
        every axis title/.append style={at={(0.1,0.8)},font=\bfseries}
	    ]
		\nextgroupplot[
		legend style = {
		  font=\small,
          draw=none, 
          fill=none,
          column sep = 1pt, 
          /tikz/every even column/.append style={column sep=5mm},
          legend columns = -1, 
          legend to name = grouplegend},
		%
		]
		
		\addplot[thick,color=tuatara,mark=*] coordinates {
          (1, 0.2797)
          (2, 0.2869)
          (3, 0.2952)
          (4, 0.2996)
          (5, 0.3028)
          (6, 0.3047)
          (7, 0.3009)
          (8, 0.2968)
          (9, 0.2265)
        }; \addlegendentry{Beauty}
         \addplot[thick,color=free_speech_aquamarine,mark=diamond*] coordinates {
          (1, 0.3831)
          (2, 0.3962)
          (3, 0.3985)
          (4, 0.4007)
          (5, 0.3961)
          (6, 0.3911)
          (7, 0.3884)
          (8, 0.3802)
          (9, 0.3569)
        }; \addlegendentry{Steam}
        \addplot[thick,color=matisse,mark=square] coordinates {
          (1, 0.6877)
          (2, 0.6955)
          (3, 0.6937)
          (4, 0.6875)
          (5, 0.6844)
          (6, 0.6794)
          (7, 0.6745)
          (8, 0.6511)
          (9, 0.6293)
        };\addlegendentry{ML-1m}
        \addplot[thick,color=sun_shade,mark=pentagon] coordinates {
          (1, 0.6781)
          (2, 0.6879)
          (3, 0.6825)
          (4, 0.6822)
          (5, 0.6810)
          (6, 0.6791)
          (7, 0.6667)
          (8, 0.6462)
          (9, 0.6167)
        };\addlegendentry{ML-20m}
        \addplot[mark size=3.6pt, color=tuatara, mark=*,
        ] coordinates {
          (6, 0.3047)
        }; 
         \addplot[mark size=4.2pt, color=free_speech_aquamarine, mark=diamond*] coordinates {
          (4, 0.4007)
        }; 
        \addplot[line width=0.6mm, mark size=3.2pt, color=matisse,mark=square] coordinates {
          (2, 0.6955)
        };
        \addplot[line width=0.6mm, mark size=3.2pt, color=sun_shade,mark=pentagon] coordinates {
          (2, 0.6879)
        };
        
        
        \nextgroupplot[
        ylabel=\large NDCG@10]
		\addplot[thick,color=tuatara,mark=*] coordinates {
          (1, 0.1626)
          (2, 0.1693)
          (3, 0.1744)
          (4, 0.1799)
          (5, 0.1809)
          (6, 0.1832)
          (7, 0.1772)
          (8, 0.1749)
          (9, 0.1264)
        };
         \addplot[thick,color=free_speech_aquamarine,mark=diamond*] coordinates {
          (1, 0.2123)
          (2, 0.2206)
          (3, 0.2226)
          (4, 0.2241)
          (5, 0.2207)
          (6, 0.2190)
          (7, 0.2149)
          (8, 0.2096)
          (9, 0.1909)
        }; 
        \addplot[thick,color=matisse,mark=square] coordinates {
          (1, 0.4673)
          (2, 0.4759)
          (3, 0.4743)
          (4, 0.4688)
          (5, 0.4650)
          (6, 0.4587)
          (7, 0.4506)
          (8, 0.4250)
          (9, 0.3942)
        };
        \addplot[thick,color=sun_shade,mark=pentagon] coordinates {
          (1, 0.4444)
          (2, 0.4513)
          (3, 0.4487)
          (4, 0.4471)
          (5, 0.4452)
          (6, 0.4434)
          (7, 0.4292)
          (8, 0.4049)
          (9, 0.3741)
        };
        \addplot[mark size=3.6pt, color=tuatara, mark=*] coordinates {
          (6, 0.1832)
        }; 
         \addplot[mark size=4.2pt, color=free_speech_aquamarine, mark=diamond*] coordinates {
          (4, 0.2241)
        }; 
        \addplot[line width=0.6mm, mark size=3.2pt, color=matisse,mark=square] coordinates {
          (2, 0.4759)
        };
        \addplot[line width=0.6mm, mark size=3.2pt, color=sun_shade,mark=pentagon] coordinates {
          (2, 0.4513)
        };
        
	\end{groupplot}
\node at ($(group c1r1) + (120pt, 95pt)$) {\ref{grouplegend}};
\end{tikzpicture}

    \caption{Performance with different mask proportion $\rho$ on $d=64$. Bold symbols denote the best scores in each line.}
    \label{fig:mask_fig}
\end{figure}


  

As described in \S~\ref{sec:ml}, mask proportion $\rho$ is a key factor in model training, which directly affects the loss function (Equation~\ref{eq:loss}).
Obviously, mask proportion $\rho$ should not be too small or it is not enough to learn a strong model.
Meanwhile, it should not be too large, otherwise, it would be hard to train since there are too many items to guess based on a few contexts in such case.
To examine this, we study how mask proportion $\rho$ affects the recommendation performances on different datasets.

Figure~\ref{fig:mask_fig} shows the results with varying mask proportion $\rho$ from 0.1 to 0.9.
Considering the results with $\rho > 0.6$ on all datasets, a general pattern emerges, the performances decreasing as $\rho$ increases.
From the results of the first two columns, it is easy to see that $\rho=0.2$ performs better than $\rho=0.1$ on all datasets.
These results verify what we claimed above.


In addition, we observe that the optimal $\rho$ is highly dependent on the sequence length of the dataset.
For the datasets with short sequence length (\textit{e.g.}, Beauty and Steam), the best performances are achieved at $\rho{=}0.6$ (Beauty) and $\rho{=}0.4$ (Steam), while the datasets with long sequence length (\textit{e.g.}, ML-1m and ML-20m) prefer a small $\rho{=}0.2$.
This is reasonable since, compared with short sequence datasets, a large $\rho$ in long sequence datasets means much more items that need to be predicted. %
Take ML-1m and Beauty as example, $\rho{=}0.6$ means we need to predict $98 {=} \lfloor 163.5{\times} 0.6 \rfloor $ items on average per sequence for ML-1m, while it is only $5 {=} \lfloor 8.8{\times} 0.6 \rfloor$ items for Beauty.
The former is too hard for model training.







\subsection{Impact of Maximum Sequence Length $N$}

\begin{table}[t]
\renewcommand{\arraystretch}{1}
    \centering
    \caption{Performance with different maximum length $N$.}
    \label{tab:seq}
    \begin{adjustbox}{max width=\linewidth}
        \begin{tabular}{l l r r r r r }
        \toprule
         & & \multicolumn{1}{c}{10} & \multicolumn{1}{c}{20} & \multicolumn{1}{c}{30} & \multicolumn{1}{c}{40}  & \multicolumn{1}{c}{50}  \\ \midrule
       \multirow{3}{*}{Beauty} & \#samples/s & 5504 & 3256 & 2284 &  1776 & 1441\\
       & HR@10 & 0.3006 & 0.3061 & 0.3057 & 0.3054 & 0.3047 \\
       & NDCG@10 & 0.1826 & 0.1875 & 0.1837 & 0.1833 & 0.1832 \\
        \bottomrule \toprule
       &  & \multicolumn{1}{c}{10} & \multicolumn{1}{c}{50} & \multicolumn{1}{c}{100} & \multicolumn{1}{c}{200}  & \multicolumn{1}{c}{400}  \\ \midrule
       \multirow{3}{*}{ML-1m} & \#samples/s & 14255 & 8890 & 5711 & 2918 & 1213\\
       & HR@10 & 0.6788 & 0.6854 & 0.6947 & 0.6955 & 0.6898 \\
       & NDCG@10 & 0.4631 & 0.4743 & 0.4758 & 0.4759 & 0.4715 \\
        \bottomrule
    \end{tabular}
    \end{adjustbox}
\end{table}




We also investigate the effect of the maximum sequence length $N$ on model's recommendation performances and efficiency.

Table~\ref{tab:seq} shows recommendation performances and training speed with different maximum length $N$ on Beauty and ML-1m.
We observe that the proper maximum length $N$ is also highly dependent on the average sequence length of the dataset.
Beauty prefers a smaller $N=20$, while ML-1m achieves the best performances on $N=200$.
This indicates that a user's behavior is affected by more recent items on short sequence datasets and less recent items for long sequence datasets.
The model does not consistently benefit from a larger $N$ since a larger $N$ tends to introduce both extra information and more noise.
However, our model performs very stably as the length $N$ becomes larger.
This indicates that our model can attend to the informative items from the noisy historical records.

A scalability concern about BERT4Rec is that its computational complexity per layer is $\mathcal{O}(n^2d)$, quadratic with the length $n$.
Fortunately, the results in Table~\ref{tab:seq} shows that the self-attention layer can be effectively parallelized using GPUs.








\subsection{Ablation Study}


Finally, we perform ablation experiments over a number of key components of BERT4Rec in order to better understand their impacts, including positional embedding (PE), position-wise feed-forward network (PFFN), layer normalization (LN), residual connection (RC), dropout, the layer number $L$ of self-attention, and the number of heads $h$ in multi-head attention. %
Table~\ref{tab:ablation} shows the results of our default version ($L=2, h=2$) and its eleven variants on all four datasets with dimensionality $d=64$ while keeping other hyper-parameters (\textit{e.g.}, $\rho$) at their optimal settings. %




We introduce the variants and analyze their effects respectively:
\begin{enumerate}%
  \item \textbf{PE}. The results show that removing positional embeddings causes BERT4Rec's performances decreasing dramatically on long sequence datasets (\textit{i.e.}, ML-1m and ML-20m). Without the positional embeddings, the hidden representation $\bm{H}_i^L$ for each item $v_i$ depends only on item embeddings.
  In this situation, we predict different target items using the same hidden representation of ``\texttt{[mask]}''. This makes the model ill-posed. This issue is more serious on long sequence datasets since they have more masked items to predict.
  \item \textbf{PFFN}. The results show that long sequence datasets (\textit{e.g.}, ML-20m) benefit more from PFFN. 
  This is reasonable since a purpose of PFFN is to integrate information from many heads which are preferred by long sequence datasets as discussed in the analysis about head number $h$ in ablation study (\ref{itm:head}). %
  \item \textbf{LN}, \textbf{RC}, and \textbf{Dropout}. These components are introduced mainly to alleviate overfitting. 
  Obviously, they are more effective on small datasets like Beauty.
   To verify their effectiveness on large datasets, we conduct an experiment on ML-20m with layer $L{=}4$.
   The results show that NDCG@10 decreases about 10\% w/o RC.
  \item \textbf{Number of layers $L$}. The results show that stacking Transformer layer can boost performances especially on large datasets (\textit{e.g}, ML-20m).
  This verifies that it is helpful to learn more complex item transition patterns via deep self-attention architecture.
  The decline in Beauty with $L=4$ is largely due to overfitting.
  \item\label{itm:head} \textbf{Head number $h$}. We observe that long sequence datasets (\textit{e.g.}, ML-20m) benefit from a larger $h$ while short sequence datasets (\textit{e.g.}, Beauty) prefer a smaller $h$. This phenomenon is consistent with the empirical result in~\cite{Tang:EMNLP2018:Why} that large $h$ is essential for capturing long distance dependencies with multi-head self-attention.
\end{enumerate}
\begin{table}[t]
\renewcommand{\arraystretch}{1}
\setlength{\tabcolsep}{0.75em}
    \centering
    \caption{Ablation analysis (NDCG@10) on four datasets. Bold score indicates performance better than the default version, while $\downarrow$ indicates performance drop more than 10\%.}
    \label{tab:ablation}
    \begin{adjustbox}{max width=\linewidth}
        \begin{tabular}{l r r r r}
        \toprule
     \multirow{2}{*}{Architecture} & \multicolumn{4}{c}{Dataset} \\ 
         \cmidrule(lr){2-5}
          & \multicolumn{1}{c}{Beauty} & \multicolumn{1}{c}{Steam} & \multicolumn{1}{c}{ML-1m} & \multicolumn{1}{c}{ML-20m} \\ \midrule
        $L=2$, $h=2$  & 0.1832 & 0.2241 & 0.4759 & 0.4513 \\ \midrule
        w/o PE  & 0.1741 & 0.2060 & 0.2155$\downarrow$ & 0.2867$\downarrow$ \\
        w/o PFFN  & 0.1803 & 0.2137 & 0.4544 & 0.4296 \\ \midrule
        w/o LN  & 0.1642$\downarrow$ & 0.2058 & 0.4334 & 0.4186 \\
        w/o RC  & 0.1619$\downarrow$ & 0.2193 & 0.4643 & 0.4483 \\
        w/o Dropout  & 0.1658 & 0.2185 & 0.4553 & 0.4471 \\ \midrule
        1 layer \hspace{\fill}($L=1$) & 0.1782 & 0.2122 & 0.4412 & 0.4238 \\
        3 layers \hspace{\fill} ($L=3$) & \textbf{0.1859} & \textbf{0.2262} & \textbf{0.4864} & \textbf{0.4661}\\
        4 layers \hspace{\fill} ($L=4$) & \textbf{0.1834} & \textbf{0.2279} & \textbf{0.4898} & \textbf{0.4732}\\ \midrule
        1 head \hspace{\fill} ($h=1$) & \textbf{0.1853} & 0.2187 & 0.4568 & 0.4402 \\
        4 heads \hspace{\fill} ($h=4$) & 0.1830 & \textbf{0.2245} & \textbf{0.4770} & \textbf{0.4520} \\
        8 heads \hspace{\fill} ($h=8$) & 0.1823 & \textbf{0.2248} & 0.4743 & \textbf{0.4550} \\
        \bottomrule
    \end{tabular}
    \end{adjustbox}
\end{table}







